\section{The Judgment at Death}

Death ends the earthly time in which we can accept or reject the grace of Christ. At that moment each person receives eternal retribution in the immortal soul: heaven immediately, heaven through purification, or hell. This is the particular judgment (CCC 1021--1022). Its particularity means that this person, with this actual history of faith and works, is referred to Christ. We shall not disappear into the fate of a crowd. Neither a family's piety nor an age's corruption can supply another person's answer in place of our own.

The Gospel gives this truth a piercing immediacy. The rich man and Lazarus die, and their condition beyond death is already different (Lk 16:19--31). The penitent thief receives Christ's promise of paradise while the Cross still stands before the spectators (Lk 23:39--43). The outward manner of dying does not disclose the inward judgment. One man possessed comfort and recognition; another lay neglected; a third died under a criminal sentence. Christ sees the soul. The earthly observer sees fragments of the story and may mistake those fragments for the whole.

We should receive both the urgency and the consolation of this teaching. Urgency follows because the time for conversion has an end. Consolation follows because our life reaches the one competent to judge it. Our destiny does not depend upon whether survivors praise us adequately, understand us fully, or preserve our memory. A holy person may die unrecognized; a great wrong may remain concealed from everyone near the deathbed. Nothing is thereby concealed from Christ. He knows the final act of repentance and the last freely maintained refusal, the interior offering and the merely outward appearance of peace.

The soul which receives this judgment is our own continuing self. The Church teaches that a spiritual element endowed with consciousness and will survives death; she calls it the soul. Bodily resurrection remains future, and concerns the whole person (CDF, \emph{Letter on Certain Questions Concerning Eschatology}, 1979, nos.~2--5). There is continuity between the person who lived, the soul judged at death, and the human being raised at the last day. Death does not cancel personal responsibility or turn the person into a memory belonging only to others.

The definitiveness of judgment gives weight to ordinary choices. During life we can repent of evil, fall from good, grow in charity, or become more deeply entangled in vice. At death this earthly probation ends. Thomas expressly joins the completion of a life to the possibility of judging it as a whole (III, q.~59, a.~5). Until its end, the history can still contain a conversion. A sinner need not continue as he has begun. A person who has served God faithfully must still persevere. The teaching directs the living toward the response that remains possible, without pretending that tomorrow belongs to us.

Postponement therefore concerns more than a date in the calendar. To say that one will repent later can be a way of choosing the present sin again. The promise of a future confession may quiet the conscience while leaving the will unchanged. Each repetition makes the familiar evil easier to protect, and each refusal of grace can make the next refusal less noticeable. God's patience gives room for repentance; we misuse it when we make it part of a plan for delaying repentance. The uncertainty of our last day strips that plan of its imagined security.

Yet the finality of death does not authorize despair over years badly spent. The life still open before Christ remains open to his mercy. The thief could not return to the victims of every past wrong, begin a long course of penance, or demonstrate years of amendment. He turned toward the Savior beside him. His conversion teaches the greatness of mercy without making delay wise. We may trust the Lord who received him; we may not presume upon an opportunity which has not been promised to us. The grace available in the present is sufficient reason to turn now.

St.~Alphonsus makes this practical turn in the twenty-fourth consideration of \emph{Preparation for Death}. After contemplating the judgment, he urges the reader to settle the account of life while repentance is possible (consideration XXIV, point~3). His point is personal and immediate. The future judgment should not remain material for religious interest while an acknowledged sin is left untouched. The sinner can examine himself, confess, accept penance, and begin to amend. The fear awakened by judgment is meant to move toward these acts, where Christ's mercy reaches an actual person with an actual need.

There is also a daily freedom in remembering the particular judgment. The opinions surrounding us can neither save nor finally condemn us. We may therefore refuse a dishonest advantage even when everyone expects us to take it. We may accept obscurity when obedience brings no praise. We may acknowledge a fault without defending an entire reputation. The measure of the life is being held by Christ. To live before him is to begin escaping the need to justify ourselves before every passing tribunal, while becoming more willing to give the honest account that charity and justice require.

